\section{Conclusions}\label{sec:conclusions}

A continuation over the weight of a normalized loss computes, for one array of point traps, a front of efficiency against uniformity of about \nLambda{} points in under a minute, with a closed-form gradient and a warm start along the weights. On a ten by ten lattice of 100 traps its hypervolume is \hvFront{} against \hvPub{} for the union of the points of Gerchberg--Saxton, the adaptive additive iteration, weighted GS and direct search, a paired difference of \hvDiff{} (95 percent interval \hvDiffLo{} to \hvDiffHi) over \nSeeds{} seeds, and it exceeds the union of the one-parameter fronts of the published methods by \rivUnionDiff{}.
The front dominates the points of Gerchberg--Saxton and of the adaptive additive iteration in every seed, and it adds to the published points the efficient side of the trade-off: the efficiency rises by at most \eMaxRatioPct{} percent over the start while the uniformity falls from one to near zero, so the cost of light in uniformity is steep. Proposition~\ref{prop:shift} gives the weight its reading, the squared error against a uniformly raised target, and the weight $\lambda^\ast=\lamSat{}$ beyond which the target exceeds the efficiency that is reached;
Proposition~\ref{prop:bound} ties the spread and the uniformity to the loss and holds on every front point.

The result is partial. The weighted GS point and the phase-fixed point are dominated by the front in \domGSWSeeds{} and \domPFSeeds{} of \nSeeds{} seeds, \frontDomMin{} to \frontDomMax{} of the points of the front are dominated by iterates of weighted GS, and the stored fronts of the benchmark of the studio are below the published union at 100 iterations on all \benchNHundred{} geometries, with a deficit of at most \benchWorstAbsHundred{}. The ablation shows that the warm start, the scaling of the loss and a start close to uniform carry the front (each raises the hypervolume by between \ablStartGsAbs{} and \ablRawAbs{} against its removal), that the steps per weight and the weights above 20 are on a plateau, and that a start from the converged hologram of weighted GS removes the domination without changing the area.

The next steps follow from these limits: start the front from the converged hologram of weighted GS; replace the knee rule by a choice of the user; give the stored fronts the weights that they need in the benchmark; put the device model of the literature in the loss, so that the front is computed for the device that displays the hologram and not for the ideal one; and measure the front on an instrument, where the trade-off between the light in the traps and their uniformity decides how many atoms an array can load.
